\subsection{Radial period coordinates and the compact quotient}
\label{mapsIII:carrier-radial-prerequisites}

The carrier constructions use the following real quotient of the period
family of \cref{wb:prop:period-input-audit}.  Fix the lattice basis
\[
 \Lambda=\mathbb Z\langle\widehat\gamma,\widehat u,
                         \widehat w,\widehat\delta\rangle.
\]
The cusp coordinate and its holomorphic residual functions give
\[
 \Pi=[Z\mid I_2],\qquad
 Z=s\begin{pmatrix}0&1\\-1&0\end{pmatrix}
       +\begin{pmatrix}6\mu(t)&h(t)\\b(t)-h(t)&\mu(t)\end{pmatrix},
 \qquad t=e^{2\pi i s}.
\]
Here $\mu,b,h$ are holomorphic on a disc containing zero.  On the positive
radial ray put $t=e^{-2\pi T}$ and $s=iT$.  Choose $T_0$ large enough that
this ray lies in that disc for $T\geq T_0$, and set
\begin{equation}
 \begin{aligned}
 \mu_T&=\alpha_T+im_T,& h_T&=\eta_T+iH_T,& b_T&=\xi_T+iB_T,\\
 c_T&=\xi_T-\eta_T,& L_T&=T+H_T,& q_T&=L_T-B_T,\\
 D_T&=L_Tq_T+6m_T^2.
 \end{aligned}
 \label{wb:cusp-spec-radial-notation}
\end{equation}
All these quantities are real.  Write
$m=\operatorname{Im}\mu(0)$, $H=\operatorname{Im}h(0)$ and
$B=\operatorname{Im}b(0)$.  Taylor expansion at zero yields
\begin{equation}
 m_T=m+O(e^{-2\pi T}),\qquad
 H_T=H+O(e^{-2\pi T}),\qquad
 B_T=B+O(e^{-2\pi T}),
 \label{wb:cusp-spec-holomorphic-errors}
\end{equation}
and the analogous estimates for the real parts.  In particular
$D_T=T^2+O(T)>0$ after increasing $T_0$.

Identify $\mathbb C^2$ with $\mathbb R^4$ in the order
$(x_1,y_1,x_2,y_2)$ and equip it with the Euclidean metric.  The induced
flat metric on a quotient is the metric used below.  This metric is an
explicit choice in these coordinates.

\begin{proposition}[Real periods and their dual coordinates]
\label{mapsIII:carrier-real-periods}
The radial fibre and its real period matrix are
\[
 F_T=\mathbb R^4/P_T\mathbb Z^4,\qquad
 P_T=\begin{pmatrix}
 6\alpha_T&\eta_T&1&0\\
 6m_T&L_T&0&0\\
 c_T&\alpha_T&0&1\\
 -q_T&m_T&0&0
 \end{pmatrix}.
\]
Its determinant is $-D_T$.  Its complete inverse transpose is
\begin{equation}
 P_T^{-\mathsf T}=
 \begin{pmatrix}
 0&0&1&0\\[2pt]
 \dfrac{m_T}{D_T}&\dfrac{q_T}{D_T}&
 \dfrac{-6m_T\alpha_T-q_T\eta_T}{D_T}&
 \dfrac{-m_Tc_T-q_T\alpha_T}{D_T}\\[7pt]
 0&0&0&1\\[2pt]
 -\dfrac{L_T}{D_T}&\dfrac{6m_T}{D_T}&
 \dfrac{6(L_T\alpha_T-m_T\eta_T)}{D_T}&
 \dfrac{L_Tc_T-6m_T\alpha_T}{D_T}
 \end{pmatrix}.
 \label{wb:cusp-spec-full-inverse-transpose}
\end{equation}
Thus $(P_T\mathbb Z^4)^*=P_T^{-\mathsf T}\mathbb Z^4$.
\end{proposition}

\begin{proof}
At $s=iT$ the complex period matrix is
\[
 Z_T=\begin{pmatrix}6\mu_T&h_T+iT\\
                          b_T-h_T-iT&\mu_T\end{pmatrix}.
\]
Taking real and imaginary parts of its columns and of the identity
columns gives the displayed $P_T$.  Expansion along its last two
columns gives $\det P_T=-L_Tq_T-6m_T^2=-D_T$.

To solve $P_T^{\mathsf T}x=(r,s,n_1,n_2)^{\mathsf T}$, its last two
equations first give $x_1=n_1$ and $x_3=n_2$.  Put
\[
 R=r-6\alpha_Tn_1-c_Tn_2,\qquad
 S=s-\eta_Tn_1-\alpha_Tn_2.
\]
The remaining two equations and their solution are
\[
 \begin{pmatrix}6m_T&-q_T\\L_T&m_T\end{pmatrix}
 \binom{x_2}{x_4}=\binom R S,
 \qquad
 \binom{x_2}{x_4}=
 \frac1{D_T}\begin{pmatrix}m_T&q_T\\-L_T&6m_T\end{pmatrix}
 \binom R S.
\]
Substituting $R,S$ proves every entry in
\eqref{wb:cusp-spec-full-inverse-transpose}.  Finally a real covector
$\ell$ is integral on $P_T\mathbb Z^4$ exactly when
$P_T^{\mathsf T}\ell\in\mathbb Z^4$, proving the dual-lattice formula.
\end{proof}

Let $e_1,\ldots,e_4$ denote the coordinate unit vectors in this real
order and let $v_{\gamma,T},v_{u,T}$ be the first two columns of $P_T$.
Then
\[
 P_T\mathbb Z^4=
 \mathbb Zv_{\gamma,T}\oplus\mathbb Zv_{u,T}
                 \oplus\mathbb Ze_1\oplus\mathbb Ze_3.
\]

\begin{proposition}[Exact compact quotient of the native cusp fibre]
\label{wb:cusp-spec-native-quotient}
Define the orthogonal projection and its projected period matrix by
\[
 \rho(x_1,x_2,x_3,x_4)=(x_2,x_4),\qquad
 Y_T=\begin{pmatrix}6m_T&L_T\\-q_T&m_T\end{pmatrix}.
\]
Then $\det Y_T=D_T>0$, so $Q_T=\mathbb R^2/Y_T\mathbb Z^2$ is a
compact real torus, and projection induces an exact sequence of compact
real Lie groups
\[
 0\longrightarrow
 \mathbb T^2_{\mathrm{re}}=
 \frac{\mathbb Re_1\oplus\mathbb Re_3}
      {\mathbb Ze_1\oplus\mathbb Ze_3}
 \longrightarrow F_T\xrightarrow{\ \pi_T\ }Q_T
 \longrightarrow0.
\]
This is a Riemannian quotient for the displayed flat metrics.  There is a
constant $C$, independent of $T\geq T_0$, such that
\[
 \|Y_T(a,b)^{\mathsf T}\|\geq(T-C)\sqrt{a^2+b^2},\qquad
 \operatorname{sys}(Q_T)\geq T-C.
\]
For every $R>0$, the radius-$R$ balls of $Q_T$ and $Q_T^2$ are exactly
Euclidean once $T-C>4R$.  Consequently their pointed limits are
$\mathbb R^2$ and $\mathbb R^4$, respectively, both locally isometrically
and smoothly, and in pointed Gromov--Hausdorff topology.
\end{proposition}

\begin{proof}
The determinant calculation gives
$\det Y_T=6m_T^2+L_Tq_T=D_T$.  Projection kills the last two period
columns and sends the first two to those of $Y_T$.  Therefore
$\rho(P_T\mathbb Z^4)=Y_T\mathbb Z^2$, which proves that $\pi_T$ is
well defined and surjective.

If $\pi_T([x])=0$, choose $a,b\in\mathbb Z$ such that
$\rho(x)=Y_T(a,b)^{\mathsf T}$.  The vector
$x-av_{\gamma,T}-bv_{u,T}$ belongs to
$\mathbb Re_1\oplus\mathbb Re_3$ and represents the same point of
$F_T$.  Conversely every point in this real plane maps to zero.
If a lattice vector lies in that plane, its first two lattice
coefficients satisfy $Y_T(a,b)^{\mathsf T}=0$.  Invertibility forces
$a=b=0$, so the intersection lattice is exactly
$\mathbb Ze_1\oplus\mathbb Ze_3$.  This proves the asserted kernel and
its injective inclusion.  Since $\rho$ is the orthogonal projection
along that plane, its Euclidean quotient metric is the displayed one.

Write
\[
 Y_T=T\begin{pmatrix}0&1\\-1&0\end{pmatrix}
       +\begin{pmatrix}6m_T&H_T\\B_T-H_T&m_T\end{pmatrix}.
\]
The second matrix has uniformly bounded operator norm by holomorphy.
Choose $C$ above this bound.  The first matrix is $T$ times an orthogonal
matrix, and the reverse triangle inequality gives the stated estimate.
A nonzero integer vector has norm at least one, which proves the
systolic bound; it also proves the same bound for the product lattice.

For $x,y$ in the Euclidean radius-$R$ ball and a nonzero lattice vector
$\lambda$, the estimate
\[
 \|x-y+\lambda\|\geq\|\lambda\|-\|x-y\|
 \geq(T-C)-2R>2R\geq\|x-y\|
\]
holds when $T-C>4R$.  Thus no lattice translation shortens a distance
between points of this ball.  Every point of the quotient ball has a
lift in the Euclidean ball, by choosing a shortest representative of
its lattice coset.  The covering map therefore gives an isometry onto
the entire quotient ball.  These are smooth local isometries, proving
all the stated pointed convergences.
\end{proof}
